%!TEX root = ../kr.tex
\chapter{Словник спеціалізованих термінів}
\label{apdx:glossary}

\noindent\textit{Нижче наведено визначення ключових термінів, що використовуються у контексті кіберзлочинності та цифрової безпеки.}

\vspace{0.5cm}

\parheading{API}~--- спосіб взаємодії між різними програмами та сервісами.

\parheading{APK}~--- файл установки додатка для Android.

\parheading{Autentification (Autentykacia)}~--- процес перевірки особи або системи для підтвердження їхньої достовірності (пароль, токен, біометрія).

\parheading{Adware}~--- програма, що відображає небажану рекламу та може стежити за користувачем.

\parheading{Blockchain}~--- технологія розподіленого реєстру, на якій побудовані крипто та NFT.

\parheading{BEC (Business Email Compromise)}~--- шахрайство через підміну ділової електронної пошти.

\parheading{Botnet}~--- мережа заражених комп'ютерів, керована зловмисником для атак та спаму.

\parheading{Caller ID}~--- функція визначення номера того, хто дзвонить (можна підмінити через VoIP).

\parheading{Cookies (Kuki)}~--- невеликі файли, що зберігаються браузером для запам'ятовування налаштувань сайту.

\parheading{Cryptocurrency (Krypto)}~--- цифрова валюта з використанням криптографії та блокчейну (BTC, USDT).

\parheading{CVV}~--- тризначний код безпеки на звороті банківської картки.

\parheading{Drive-by Download}~--- завантаження шкідливого кода з сайту без відома користувача.

\parheading{Exploit}~--- програма або код, що використовує вразливість системи для несанкціонованого доступу.

\parheading{Formjacking}~--- атака, при якій зловмисник перехоплює дані з веб-форм.

\parheading{FOMO (Fear Of Missing Out)}~--- синдром втраченої можливості, почуття страху пропустити щось важливе (використовується в шахрайстві для маніпуляції).

\parheading{Homograph Attack (Homoglif-ataka)}~--- використання символів, схожих один на одного (a/а), для створення фальшивих веб-адрес.

\parheading{Keylogging}~--- таємне запис натискання клавіш для крадіжки паролів і персональних даних.

\parheading{KYC (Know Your Customer)}~--- процедура перевірки особи клієнта у фінансових сервісах.

\parheading{Malware (Zlovred)}~--- шкідливе програмне забезпечення (віруси, трояни, шпигунські програми, черв'яки).

\parheading{NFT}~--- невзаємозамінний токен, унікальний цифровий актив на блокчейну.

\parheading{OTP (One-Time Password)}~--- одноразовий пароль, що приходить в SMS або генерується додатком.

\parheading{Overlay (Ataka)}~--- фальшиве вікно поверх справжнього додатка або веб-сайту для крадіжки даних.

\parheading{Phishing (Fishing)}~--- розсилання фальшивих листів або посилань для крадіжки персональних даних, паролів та грошей.

\parheading{Ponzi Scheme}~--- фінансова піраміда, схема, що обіцяє високий прибуток за рахунок коштів нових учасників.

\parheading{Ransomware (Prohrama-vymahach)}~--- вірус, що блокує файли та вимагає викупу за їх розблокування.

\parheading{RAT (Remote Access Trojan)}~--- програма-шпигун, що дає повний віддалений доступ до вашого пристрою.

\parheading{Rootkit}~--- програма, що дає повний контроль над системою на глибокому рівні.

\parheading{Session (Sesiya)}~--- період активної роботи користувача на сайті або в додатку.

\parheading{Sideloading}~--- установка додатків для Android не через Google Play, часто використовується для розповсюдження шкідливого коду.

\parheading{Smishing}~--- фішинг через SMS-повідомлення з посиланнями на фальшиві сайти.

\parheading{Social Engineering (Sotsialna inzheneriya)}~--- маніпуляція людьми для отримання конфіденційної інформації або доступу до систем.

\parheading{Spoofing (Spufing)}~--- імітація справжнього відправника для обману (поддільний номер телефону, адреса електронної пошти).

\parheading{Spyware}~--- програма, що таємно стежить за діяльністю користувача та збирає персональні дані.

\parheading{Token}~--- цифровий код або криптовалюта для доступу або платежів.

\parheading{Trojan (Troyan)}~--- вірус, замаскований під корисну програму, що дає зловмиснику доступ до вашого пристрою.

\parheading{Verification (Verifikacia)}~--- перевірка автентичності користувача або об'єкта.

\parheading{Vishing (Vishin)}~--- фішинг через голосові дзвінки з числення видавання себе за доказові установи.

\parheading{Wallet (Gamanets)}~--- програма або пристрій для зберігання криптовалюти та цифрових активів.

\parheading{Worm (Chervak)}~--- самовідтворюваний вірус, що розповсюджується мережею без участі користувача.

\parheading{Zero-day}~--- вразливість у програмному забезпеченні, невідома розробнику, що дозволяє зловмисникам атакувати до виходу оновлення.

\parheading{Кібершахрайство}~--- злочинна діяльність з використанням інтернету та цифрових технологій для обману та крадіжки коштів.

\parheading{Цифровий імунітет}~--- запас знань і навичок для розпізнавання та уникнення цифрових загроз.

\parheading{Платформи оголошень}~--- веб-сайти та мобільні програми для розміщення купівлі-продажу товарів (OLX, Avito тощо).

\parheading{Фішинг}~--- атака через фальшиві листи та посилання для крадіжки даних.

\parheading{Вешинг}~--- атака через голосові дзвінки з выдачею себе за компанії.

\parheading{Соціальна інженерія}~--- маніпуляція людьми для отримання конфіденційної інформації.

\parheading{Фейкові акаунти}~--- фальшиві профілі в соціальних мережах, створені для шахрайства та маніпуляції.

\parheading{Маркетплейси}~--- великі онлайн-платформи для замовлення товарів та послуг (Amazon, eBay, Wildberries).

\parheading{Передоплата}~--- запит грошей перед доставкою товару (часто використовується в шахрайстве).

\parheading{Посередник}~--- третя сторона, на яку припускають під час трансакції для збереження коштів (практика, яка часто використовується в шахрайстві).

\parheading{Дропер}~--- шахрайський посередник, що отримує та передає крадені гроші.

\parheading{Мула}~--- особа, яку використовують для передачі крадених коштів через свій банківський рахунок.

\parheading{Маніпуляція}~--- вплив на людину для досягнення певної поведінки на користь маніпулятора.

\parheading{Терміновість}~--- створення ілюзії нехватки часу для прийняття поспішних рішень.

\parheading{Страх}~--- психологічна емоція, що використовується для маніпуляції поведінкою.

\parheading{Довіра}~--- сприйняття людини як надійної особи чи організації.

\parheading{Емпатія}~--- здатність розуміти та розділяти почуття іншої людини (використовується в соціальній інженерії).

\parheading{Цинізм}~--- недовіра до людей та їхніх намірів.

\parheading{Дефіцит часу}~--- обмеженість часу для прийняття рішення, що впливає на якість аналізу інформації.

\parheading{Верифікація}~--- процедура для підтвердження особи власника документів.

\parheading{Вішинг (Vishing)}~--- телефонне шахрайство, коли зловмисники видають себе за банк або поліцію.

\parheading{Двофакторна автентифікація (2FA)}~--- подвійний захист: пароль + код.

\parheading{Демпінг}~--- від англ. dumping. Продаж товарів за цінами, нижчими від собівартості.

\parheading{Домен}~--- адреса сайту в інтернеті (наприклад, google.com).

\parheading{Зловред (Malware)}~--- загальний термін для шкідливого програмного забезпечення (віруси, трояни, черв'яки, шпигуни).

\parheading{Імперсонація}~--- видавання себе за іншу людину, підробка особистості в інтернеті.

\parheading{Інтерфейс}~--- зовнішній вигляд сайту чи програми.

\parheading{ІПН}~--- реєстраційний номер облікової картки платника податків.

\parheading{Кейлоггер (Keylogger)}~--- програма для прихованого запису натискань клавіш з метою викрадення паролів.

\parheading{Клієнт-банк}~--- програма або сайт для керування банківськими рахунками онлайн.

\parheading{Колектори}~--- люди, що вимагають повернути борг.

\parheading{Криптографічний ключ}~--- секретна або публічна послідовність для шифрування та розшифрування даних.

\parheading{Макроси}~--- приховані команди в документах Word або Excel, які можуть виконувати шкідливі дії.

\parheading{Маркетплейс}~--- онлайн-платформа схожа на електронний базар або ярмарку.

\parheading{Месенджер}~--- програма в телефоні для листування та дзвінків.

\parheading{Нікнейм}~--- псевдонім або вигадане ім'я користувача в інтернеті.

\parheading{Павербанк}~--- портативний зовнішній акумулятор, який дозволяє заряджати ваші електронні пристрої.

\parheading{ПІН}~--- секретний код для доступу до картки або телефону (зазвичай 4 цифри).

\parheading{Піддомен}~--- додаткова адреса, що створюється перед основним доменом (наприклад, login.bank.ua).

\parheading{Підміна DNS (DNS Spoofing)}~--- перенаправлення користувача на фальшивий сайт шляхом зміни DNS-записів.

\parheading{Проксі}~--- проміжний сервер для анонімного доступу в інтернет.

\parheading{Реквізити}~--- дані про картку, необхідні для переказу грошей (номер, CVV, строк дії).

\parheading{Ре­зерв­не ко­пі­ю­ван­ня (Backup)}~--- створення копій даних для відновлення після втрати.

\parheading{Селфі}~--- фотографія самого себе, зроблена на камеру смартфона.

\parheading{Скріншот}~--- знімок екрана пристрою.

\parheading{Смішинг (Smishing)}~--- шахрайство через SMS-повідомлення з підозрілими покликаннями.

\parheading{Соціальна інженерія}~--- маніпулятивні техніки впливу на людей для отримання конфіденційної інформації.

\parheading{Термінал}~--- пристрій для оплати картками.

\parheading{Токен (Token)}~--- тимчасовий ключ доступу замість пароля, або спеціальний цифровий підпис для підтвердження платежів.

\parheading{Транзакція}~--- будь-яка фінансова операція (переказ, оплата, зняття грошей).

\parheading{Тригер}~--- первинний сигнал, який запускає настороженість і спонукає зупинитися та перевірити.

\parheading{Троль (Інтернет-троль)}~--- користувач, який навмисно провокує конфлікти та поширює суперечливий контент в інтернеті.

\parheading{Троян (Trojan)}~--- шкідлива програма, що маскується під корисний додаток.

\parheading{Троянська програма}~--- зловред, що маскується під легітимне ПЗ для прихованого доступу до системи.

\parheading{Фішинг (Phishing)}~--- обман через підроблені сайти або листи для викрадення паролів і даних карток.

\parheading{Фішингове покликання}~--- підроблене покликання з метою обману.

\parheading{Хеш (Hash)}~--- унікальний цифровий «відбиток» файлу для перевірки його справжності.

\parheading{Хеш-функція}~--- одностороннє перетворення даних у фіксований рядок для перевірки цілісності.

\parheading{Чат}~--- місце для обміну повідомленнями в інтернеті (групова розмова).

\parheading{Цифровий двійник (Deepfake)}~--- відео, аудіо або зображення, створені штучним інтелектом для імітації реальної людини. Використовується шахраями для дзвінків «від родичів» або відеозвернень «від керівництва».

\parheading{Цифровий сертифікат}~--- електронний документ, що підтверджує автентичність веб-сайту або особи.

\parheading{Шифрування}~--- процес перетворення інформації у нечитабельний вигляд для захисту від несанкціонованого доступу.

\parheading{Шпигунське ПЗ (Spyware)}~--- програми для прихованого збору даних про активність користувача.
